\DiaryEntry{Book of Integrals, 1}{2026-09-18}{Integrals}

This is based on \todo{get bibtex entry for book}

Start with

\begin{equation*}
    \frac{d a^x}{dx}
\end{equation*}

We can rewrite $a = e^{\ln a}$ and therefore $a^x = e^{x \ln a}$ and this is something we can differentiate,

\begin{equation*}
    \frac{d a^x}{dx} = \frac{d e^{x \ln a} }{dx} = \ln a e^{x \ln a} = a^x \ln a \qed
\end{equation*}

Reading this from right to left, we get the integration rule,

\begin{equation*}
    \int a^x dx = \frac{1}{\ln a} a^x + C \qed
\end{equation*}


\subsubsection{Polynomials}

Rather simple, everything is based on

\begin{equation*}
    \int x^k dk = \frac{x^{k+1}}{k+1} + C \quad \forall k \neq -1
\end{equation*}

\paragraph{Example 2.} We have

\begin{equation*}
    \int 6 \sqrt{x} - \frac{4}{x^3} dx = \int 6 x^{1/2} - 4 x^{-3} dx = 6 \frac{x^{3/2}}{\frac{3}{2}} - 4 \frac{x^{-2}}{(-2)} = 4 x^{3/2} + \frac{2}{x^2}
\end{equation*}

\paragraph{Example 30.} We have

\begin{equation*}
    \int \frac{1}{1 - \sin x} \frac{\sin(x)}{1 + \sin x} dx = \int \frac{\sin x}{1 - \sin^2 x} dx = \int \frac{\sin x}{\cos^2 x} dx = \int \frac{\tan x}{\cos x} dx = \int \tan x \sec x dx = \sec x
\end{equation*}

Here, the last integral is "knowledge".

\paragraph{Example 31.} In this example, factoring the expression is the key to the solution.

\begin{equation*}
    \int \sqrt{e^{2x} - 6 e^{x} + 9} dx = \int \sqrt{(e^x - 3)^2} dx = \int e^x - 3 dx = e^x - 3x
\end{equation*}


\paragraph{Example 33.} In a similar spirit, factoring helps

\begin{equation*}
    \int \frac{6x^2 + 12x + 6}{x (x+1)^2} dx = 6 \int \frac{x^2 + 2x + 1}{x (x+1)^2} dx = 6 \int \frac{(x+1)^2}{x (x+1)^2} dx = \cdots
\end{equation*}

This greatly simplifies the integrand and we obtain

\begin{equation*}
    \cdots = \int \frac{6}{X} dx = 6 \ln |x|
\end{equation*}

\subsubsection{Substitution}

The trick is to make a "good" substitution; just be careful, that $dx$ has to be transferred into $du$.

\paragraph{Example 41.} We start with a simple integral

\begin{equation*}
    \int \frac{4}{x+8} dx
\end{equation*}

We substitute $u = x+8$, from which we have $du/dx = 1 \rightarrow dx = du$ and the integral becomes

\begin{equation*}
    \int \frac{4}{u} du = 4 \ln|u| = 4 \ln |x+8|
\end{equation*}

The importan thing to notice is that we divide by the derivative of the substitution; this may prevent the substitution from working as in the following example.

\begin{equation*}
    \int \sin x^2 dx
\end{equation*}

Setting $u = x^2$, we have $du/dx = 2x \rightarrow dx = du / 2x$ and the integal becomes

\begin{equation*}
    \int \sin u \frac{du}{2x} = ?
\end{equation*}

This integral contains both $u$ and $x$ and therefore cannot be solved. We have hit a dead end. Checking the integral, it turns out that it cannot be expressed in closed-form.

The following integral, on the other han, is solvable by the same substitution,

\begin{equation*}
    \int x \sin(x^2) dx = \int x \sin u \frac{du}{2x} = \frac{1}{2} \int \sin u du = - \frac{1}{2} \cos u = - \frac{1}{2} \cos x^2 \qed
\end{equation*}

While we are at it, we can calculate the definitve integral

\begin{equation*}
    \int_0^x t \sin t dt = - \left. \frac{1}{2}  \cos t^2 \right|_0^x = \frac{1}{2} (1 - \cos x^2) \qed
\end{equation*}

\todo{plot}


\paragraph{Example 42.} We consider

\begin{equation*}
    \int x (3x^2 + 1)^4 dx
\end{equation*}

Using $u = 3x^2 + 1$, we have $du/dx = 6x \rightarrow dx = du / 6x$ and we have

\begin{equation*}
    \int x u^4 \frac{du}{6x} = \frac{1}{6} \int u^4 du = \frac{u^5}{30} = \frac{(3x^2+1)^5}{30}
\end{equation*}

\paragraph{Example 58.} When we have roots, the "usual" trick is to substitute the expression inside the root. We have

\begin{equation*}
    \int 4 x^5 \sqrt{x^3+1} dx
\end{equation*}

Setting $u = \sqrt{x^3+1}$ gives $u^2 = x^3 + 1$ which we can differentiate to obtain $2u du = 3x^2 dx \rightarrow dx = \frac{2}{3} \frac{u}{x^2} du$ which allows us to rewrite the integral as
\begin{equation*}
    \int 4x^5 u \frac{2}{3} \frac{u}{x^2} du = \frac{8}{3} \int x^3 u^2 du 
\end{equation*}

The trick is to express $x^3$ by means of $u^2 = x^3 + 1$: We have $x^3 = u^2 - 1$ and the integral becomes

\begin{equation*}
    \frac{8}{3} \int (u^2 - 1) u^2 du = \cdots = \frac{8}{15} \sqrt{x^3+1}^5 - \frac{8}{9} \sqrt{x^3+1}^3 \qed
\end{equation*}

\paragraph{Example 59.} The same trick works for

\begin{equation*}
    \int x^3 \sqrt{x^2-1} dx
\end{equation*}

By setting $u = \sqrt{x^2-1}, u^2 = x^2-1$ we have $2u du = 2x dx \rightarrow dx = u/x du$. We further obtain

\begin{equation*}
    \int x^3 u \frac{u}{x} du = \int x^2 u^2 du = \int (u^2+1) u^2 du = \cdots
\end{equation*}

The other substitution option is to set $u=x^2 - 1$ from which we have $du/dx = 2x \rightarrow dx = du/2x du$ and we get

\begin{equation*}
    \int x^3 \sqrt{u} \frac{du}{2x}du = \int \frac{1}{2} x^2 u^{1/2} du = \frac{1}{2} \int (u+1)u^{1/2} du = \cdots
\end{equation*}

I don't tink that there are cases where one option substitution option actually solves the problem whereas the other option does not.

This root trick does not always work; see the following variation of Exmaple 58 from above,

\begin{equation*}
    \int x^6 \sqrt{x^3+1} dx = \frac{2}{3} \int x^6 u^2 \frac{du}{x^2} = \frac{2}{3} \int x^4 u^2 du 
\end{equation*}

Continuiing as above, we have $x^3 = u^2 - 1$ and this yields

\begin{equation*}
    \frac{2}{3} \int x \cdot x^3 u^2 du = \frac{2}{3} \int x \cdot (u^2 - 1) u^2 du
\end{equation*}

But now we have hit a dead end as we can't remove the $x$-term. The integral needs to be solved wit hother mena s\todo{which ones?}.

\paragraph{Exmaple 63.} This one is looks rathe spectacular,

\begin{equation*}
    \int \frac{x^5}{\sqrt{x^2 -1}} dx
\end{equation*}

However, setting $u=x^2-1$, we obtain $du/dx = 2x \rightarrow dx = du/(2x)$ and we continue as

\begin{equation*}
    \int \frac{x^5}{\sqrt{u}} \frac{du}{2x} = \frac{1}{2} \int x^4 u^{-1/2} du = \frac{1}{2} \int (u+1)^2 u^{-1/2} du = \cdots \qed
\end{equation*}

where we have again used $x^2 = u+1 \rightarrow x^4 = (u+1)^2$.

\paragraph{Further Example.} We can play around with integrals of the form

\begin{equation*}
    \int x^k \sqrt{1+x^2} dx
\end{equation*}

Using the substitution $u=1+x^2$, we have $du/dx = 2x \rightarrow dx = du/(2x)$ and we get

\begin{equation*}
    \int x^k u^{1/2} \frac{du}{2x} = \frac{1}{2} \int x^{k-1} u^{1/2} du  
\end{equation*}

From this we see that the integral is solvable for $k=1$; ie $\int x \sqrt{1+x^2} dx$. Expressing $x^2 = u-1$, we see that $k-1$ has to be even so that $x^{k-1}$ can be epxressed in terms of $x^2 = u-1$. This implies that the integral $\int x^k \sqrt{1+x^2} dx$ becomes solvable for $k$ being odd. \qed

%%% Local Variables:
%%% mode: latex
%%% TeX-master: "journal"
%%% End:
